\documentclass{article}

\usepackage{fullpage}
\usepackage{url}
\usepackage{graphicx}
\usepackage{amsmath}
\usepackage{physics}
\DeclareMathOperator{\erfi}{erfi}

\title{Quantum Mechanics II Excercise 1}
\author{Idan Stark}
\date{\today}

\begin{document}
\maketitle
\section{Natrual Units}
\label{sec:ex1}
\subsection{Unit Equivilence}
In CGS, the speed of light is measured in $cm/sec$. Therefore, for a unit system where $c = 1$, we get:
\begin{equation}
    c = 3\cdot10^{10}\:cm / sec = 1
\end{equation}
\begin{equation}
    1 sec = 3\cdot10^{10}\:cm
\end{equation}
Therefore, distance and time are measured using the same units. Using the same logic, in CGS, mass is measured in grams and energy is measured in ergs. Assigning equation (2) into the definition of 1 erg, we get:
\begin{equation}
    1\:erg = 1\:g\cdot\frac{cm^2}{sec^2} = 1\:g\cdot\frac{cm^2}{9\cdot10^{20}\:cm^2} = \frac{1}{9}\cdot10^{-20}\:g
\end{equation}
And so we got that, assuming $c = 1$, the units for mass and energy are unified. This can also be shown using the relation from special relativity $E = mc^2$. As instructed, the units of mass and energy will be $eV$.
\newline
Now let's use the fact that $\hbar = 1$. $\hbar$ Has units of angualr momentum. Angular momentum has units of $energy \times time$. Since $\hbar$ is unitless, this means that angular momentum is unitless, and requires the units of energy to be opposite to the unit of time:
\begin{equation}
    \hbar = 10^{-27}\:erg\cdot sec = 1
\end{equation}
\begin{equation}
    1\:sec = 10^{27}\:erg^{-1}
\end{equation}
Therefore the units of distance and time are opposite to the unit of energy and mass, and $eV^{-1}$ can be used to measure them. Writing all of the units in ${eV}$ is now just an issue of conversion from erg. 
\begin{equation}
    1\:erg = 6.242\cdot10^{11}\:eV
\end{equation}
\begin{equation}
    1\:g = 9\cdot10^{20}\:erg = 5.617\cdot10^{32}\:eV
\end{equation}
\begin{equation}
    1\:sec = 10^{27}\:erg^{-1} = 1.6\cdot10^{15}\:eV^{-1}
\end{equation}
\begin{equation}
    1\:cm = sec/(3\cdot10^{10}) = 5.34\cdot10^{4}\:eV^{-1}
\end{equation}
\subsection{Newton's Constant and Plank Mass}
Using the values from equations (6) to (9):
\begin{equation}
    G_N = 6.7\cdot10^{-8} cm^3sec^{-2}g^{-1}
\end{equation}
\begin{equation}
    G_N = 6.7\cdot10^{-8}\cdot(5.34\cdot10^{4}\:eV^{-1})^{3}\cdot(1.6\cdot10^{15}\:eV^{-1})^{-2}\cdot(5.617\cdot10^{32}\:eV)^{-1}
\end{equation}
\begin{equation}
    G_N = 7.095\cdot10^{-57}\:eV^{-2}
\end{equation}
This is the very close to the value we would get using Plank Mass $M_{pl} = 1.2\cdot10^{19}\:GeV$:
\begin{equation}
    G_N = M_{pl}^{-2} = (1.2\cdot10^{19}\:GeV)^{-2} = 6.94\cdot10^{-57}\:eV^{-2}
\end{equation}
With a difference of around $2.1\%$, resulting probably from our rounding of the numbers.
\subsection{The Actions of different systems}
\subsubsection{Tennis ball in a court}
The average serve velocity is around $174 \ km/h = 48.33 \ m/s $\footnote{https://toomanyrackets.com/average-serve-speeds-in-tennis/}. The length of a tennis court is $23.78$ meters. Assuming the ball is being served horizontally, the time it would take to complete that distance is $0.492$ seconds. In that time, the ball drops
\begin{equation}
    \Delta y = \frac{1}{2}gt^2 = 1.18 \ m
\end{equation}
Assuming the pitch is done from an average height of 1.5 meters, we can ignore the floor. We will also ignore air resistance, even though it affects the motion, it will not make a visible change in the action relative to that of the water molecule. The speed of the ball is then, as a function of time:
\begin{equation}
    \dot{x}^2 = v_0^2 + (gt)^2 = 2335.79 + 96.04t^2
\end{equation}
\begin{equation}
\begin{gathered}
    \frac{1}{\hbar}S[x_c(t)] = \\
    = \frac{1}{\hbar} \int_{0}^{0.492 \ sec} \frac{1}{2} 0.056 \ kg \cdot \left( 2335.79 + 96.04t^2\right)dt \sim 3.2 \cdot 10^{35}
\end{gathered}
\end{equation}
\subsubsection{Water molecule in a teacup}
The mean free path of a water molecule is roughly $2.5$ \AA \footnote{https://galileo.phys.virginia.edu/classes/304/h2o.pdf}. During this time, the molecule moves in an average velocity depending on the temperature, following the formula
\begin{equation}
    v = \sqrt{\frac{2 k_B T}{m}}
\end{equation}
Assuming the temperature of the tea is roughly $60^\circ C$ and the mass of a water molecule is $3 \cdot 10^{-26} \ kg$, we get and average speed of $553.75 \ m/s$. This makes the mean free time $T = 4.5 \cdot 10^{-13}$ seconds. Assigning all of these into the action formule, we get:
\begin{equation}
\begin{gathered}
    \frac{1}{\hbar}S[x_c(t)] = \\
    = \frac{1}{\hbar} \int_{0}^{4.5\cdot10^{-13}} \frac{1}{2} 10^{-26} \ kg \cdot \left(553.75 \ m/s\right)^2 \ dt \sim 21
\end{gathered}
\end{equation}
\section{Smeared Amplitude}
\subsection{$K\left(\vec{x}, t\right)$ is a distribution: Intuitive explanation}
In the integration over $d^3x$, the $r$ in the denominator would cancel out with the $r$ in the Jacobian. As for the integrand, the sapid oscilations near its singularities could be canceled by a bounded test function, resulting in a finite numerical value.
\subsection{Smeared Amplitude of a 3D Gaussian}
\begin{equation}
    \overline{K}\left(\vec{X}, t\right) = \int d^3x \left(\frac{1}{\pi l^2}\right)^{3/2}e^{-\left(\frac{\vec{x} - \vec{X}}{l}\right)^2} \frac{-i}{\left(2\pi\right)^2 r}
    \int_{-\infty}^{\infty} dp p e^{i\left(pr - t\sqrt{m^2 + p^2}\right)}
\end{equation}
Without limitation of generality, we can set our coordinate system such that X is on the z axis for the sake of integration. This means that
\begin{equation}
    \left(\vec{x} - \vec{X}\right)^2 = r^2 + R^2 - 2\vec{x}\cdot\vec{X} = r^2 + R^2 - 2rR \cos(\theta)
\end{equation}
Where $R = |\vec{X}|$ and $\vec{x} = (r, \theta, \varphi)$ in spherical coordinates. This makes our integral into
\begin{equation}
\begin{gathered}
    \overline{K}\left(\vec{X}, t\right) = \int d^3x \left(\frac{1}{\pi l^2}\right)^{3/2}e^{-\frac{r^2 + R^2 - 2rR\cos(\theta)}{l^2}} \frac{-i}{\left(2\pi\right)^2 r}
    \int_{0}^{\infty} dp p e^{i\left(pr - t\sqrt{m^2 + p^2}\right)} = \\
    = \left(\frac{1}{\pi l^2}\right)^{3/2} \frac{-i}{\left(2\pi\right)^2}
    \int_{0}^{\infty} dp p e^{-it\sqrt{m^2 + p^2}}
    \int_0^{\infty} dr \int_0^{\pi} d\theta \int_0^{2\pi} d\varphi
    \ r^2 \sin(\theta)
    \frac{1}{r}
    \exp[-\frac{r^2 + R^2 - 2rR\cos(\theta)}{l^2} + ipr]
\end{gathered}
\end{equation}
The integral over $\varphi$ gives $2\pi$, while the integrals over $r$ and $\theta$ are gaussians. Starting from the integral over $r$ we perform a variable change
\begin{equation}
\begin{gathered}
    u = \frac{r - v}{l} \qquad ; \qquad v = R\cos(\theta) + \frac{ipl^2}{2} \\
    u^2 = \frac{1}{l^2}\left[
        r^2 + \left(R\cos(\theta) + \frac{ipl^2}{2}\right)^2 - 2Rr\cos(\theta)
    \right] - ipr \\
    r = lu + v \qquad ; \qquad dr = ldu
\end{gathered}
\end{equation}
Which gives the $r$ integral
\begin{equation}
\begin{gathered}
    \int_0^\infty dr r \exp[-\frac{r^2 + R^2 - 2rR\cos(\theta)}{l^2} + ipr] = \\
    = l \exp[-\frac{R^2}{l^2} + \frac{R\cos(\theta) + \frac{ipl^2}{2}}{l^2}] \int_0^\infty du (lu + v) e^{-u^2}
\end{gathered}
\end{equation}
The integral in the last part evaluates to
\begin{equation}
\begin{gathered}
    \int_0^\infty du (lu + v) e^{-u^2} = \\
    = l\int_0^\infty u e^{-u^2} du + v \int_0^\infty e^{-u^2} du = \\
    = -\frac{l}{2} \left[e^{-u^2}\right]_0^\infty + \frac{\sqrt{\pi}v}{2} = 
    \frac{1}{2}\left[\sqrt{\pi}v + l\right]
\end{gathered}
\end{equation}
Assigning it back into equation 21 using the expression from the last part of equation 23, and assigning the integral over $\varphi$ to be $2\pi$, we get:
\begin{equation}
\begin{gathered}
    \overline{K}\left(\vec{X}, t\right) =  \\
    = \left(\frac{1}{\pi l^2}\right)^{3/2} \frac{-i}{2\pi}
    \int_{0}^{\infty} dp p e^{-it\sqrt{m^2 + p^2}}
    \int_0^{\pi} d\theta
    \ \sin(\theta)
    l \exp[-\frac{R^2}{l^2} + \frac{R\cos(\theta) + \frac{ipl^2}{2}}{l^2}]
    \frac{1}{2}\left[\sqrt{\pi}\left(R\cos(\theta) + \frac{ipl^2}{2}\right) + l\right]
\end{gathered}
\end{equation}
Moveing to integral over $v$ as it is defined in equation 22, and noticing that $\sin(\theta)d\theta = dv/R$:
\begin{equation}
\begin{gathered}
    \overline{K}\left(\vec{X}, t\right) =  \\
    = \left(\frac{1}{\pi l^2}\right)^{3/2} \frac{-il}{2\pi R}
    e^{-\frac{R^2}{l^2}}
    \int_{0}^{\infty} dp p e^{-it\sqrt{m^2 + p^2}}
    \int_{v_{+}}^{v_{-}} dv
    \frac{1}{2}\left[\sqrt{\pi}v + l\right]e^{\frac{v^2}{l^2}}
\end{gathered}
\end{equation}
Where $v_{\pm} = v(\cos(\theta) = \pm 1) = \pm R + \frac{ipl^2}{2}$ Which again gives the complex gaussian integral, over the variable $iv$ this time, and with finite limts. The integral of the complex gaussian function over a subsection of the plane is known as the imaginary error function $\boldsymbol{\erfi}$.
\begin{equation}
    \int_{v_{+}}^{v_{-}} dv
    \frac{1}{2}\left[\sqrt{\pi}v + l\right]e^{\frac{v^2}{l^2}} = \frac{1}{2}\left[
        \sqrt{\pi}l^2\left(e^\frac{v_{-}^2}{l^2} - e^\frac{v_{+}^2}{l^2}\right) +
        l^2 \left(\erfi\left(\frac{v_{-}}{l}\right) - \erfi\left(\frac{v_{+}}{l}\right)\right)
    \right]
\end{equation}
And therefore
\begin{equation}
\begin{gathered}
    \overline{K}\left(\vec{X}, t\right)
    = \left(\frac{1}{\pi l^2}\right)^{3/2} \frac{-il}{2\pi R}
    e^{-\frac{R^2}{l^2}}
    \int_{0}^{\infty} dp p e^{-it\sqrt{m^2 + p^2}}
    \frac{1}{2}\left[
        \sqrt{\pi}l^2\left(e^\frac{v_{-}^2}{l^2} - e^\frac{v_{+}^2}{l^2}\right) +
        l^2 \left(\erfi\left(\frac{v_{-}}{l}\right) - \erfi\left(\frac{v_{+}}{l}\right)\right)
    \right] = \\
    = \frac{1}{2\pi^2 R}
    \int_{0}^{\infty} dp p e^{-it\sqrt{m^2 + p^2}}
    \left[
        e^{-\frac{(pl)^2}{4}}\sin(pR) +
        \frac{i}{2 \sqrt{\pi}} e^{-\frac{R^2}{l^2}} \left(\erfi\left(\frac{v_{-}}{l}\right) - \erfi\left(\frac{v_{+}}{l}\right)\right)
    \right]
\end{gathered}
\end{equation}
Using the definition of $erfi = -ierf(ix)$, we get the final value
\begin{equation}
    \overline{K}\left(\vec{X}, t\right)
    = \frac{1}{2\pi^2 R}
    \int_{0}^{\infty} dp p e^{-it\sqrt{m^2 + p^2}}
    \left[
        e^{-\frac{(pl)^2}{4}}\sin(pR) +
        \frac{1}{2 \sqrt{\pi}} e^{-\frac{R^2}{l^2}}
        \left(\erf\left(\frac{R}{l} - \frac{ipl}{2}\right) - \erf\left(-\frac{R}{l} - \frac{ipl}{2}\right)\right)
    \right]
\end{equation}
If we take the limit $R \gg l$, the error functions go to $\pm 1$ at infinity, but the exponential factor in their terms goes to zero, and so that whole term disappears and we are left with
\begin{equation}
    \overline{K}\left(\vec{X}, t\right)_{R \gg l}
    = \frac{1}{2\pi^2 R}
    \int_{0}^{\infty} dp p e^{-it\sqrt{m^2 + p^2}} e^{-\frac{(pl)^2}{4}}\sin(pR)
\end{equation}
Opening up the sine will give the same expression as seen in the excercise, but twice, with opposite signs. Switching the sign of the second term will give us the same expression as in the excercise.
\begin{equation}
    \overline{K}\left(\vec{X}, t\right)_{R \gg l}
    = \frac{1}{2\pi^2 R}
    \int_{0}^{\infty} dp p e^{-\frac{(pl)^2}{4}} e^{i\left(pR - t\sqrt{m^2 + p^2}\right)}
\end{equation}

\subsection{Closed-Form Extimate of $\overline{K}\left(\vec{X}, t\right)$}
To integrate over $p$, we must use a contour Couchy integration without passing through areas where the function is singular. The damping term $e^{-\frac{(pl)^2}{4}}$ diverges when $-\Re{p^2} = |p|^2 (\sin^2(\arg(p)) - \cos^2(\arg(p))) > 0$ meaning which happens in the top and bottom quarters of the plane $|\cos(\arg(p))| > |\sin(\arg(p))|$. This, of course, is only a problem if $p$ goes to infinity in this region. If we limit it to go to infinity only on the real line, the problem would vanish.
\newline
Meanwhile, using the approximation $R \gg t$, the saddle point is $p = \pm im$. In order to work with both of these limitations, we will choose a countour that starts from $-\infty$, goes to $-\varepsilon$ where $\varepsilon\rightarrow0$ and $\varepsilon>0$, then goes up along the imaginary axis with values $-\varepsilon+iy$ until it passes through the saddle point, then goes down along the line $+\varepsilon + iy$ until it reachs the real line, and from there to infinity.
This still leaves the problem of the singularity of the exponent at the saddle point. To avoid this problem we will use a Taylor approximation to second order $pR - t\sqrt{m^2 + p^2} \approx pr - \frac{it}{2}\left(m^2 + p^2\right)$.
Now we are ready to use the saddle point approximation, we start by calculating the second derivative of the exponent at the point and taking its argument $\theta = arg (-it) = 1.5\pi$. This gives the value of the integral
\begin{equation}
    I_C = im e^{\frac{l^2m^2}{4}}e^{-i\pi/4}\sqrt{\frac{2\pi}{Rt}} = e^{i\pi/4} m e^{\frac{l^2m^2}{4}}\sqrt{\frac{2\pi}{Rt}}
\end{equation} 
\section{Discrete version of the free scalar field}
\subsection{Commutation relations based on the Fourier modes}
\begin{equation}
\begin{gathered}
    \left[\phi_n, \phi_m\right]
    = \left[\int_{-\pi}^{+\pi} \frac{dk}{2\pi}e^{ink}\phi(k), \int_{-\pi}^{+\pi} \frac{dp}{2\pi}e^{imp}\phi(p)\right]
    = \int_{-\pi}^{+\pi} \frac{dk}{2\pi} \int_{-\pi}^{+\pi}  \frac{dp}{2\pi} e^{i(nk + mp)} \left[\phi(k), \phi(p)\right] = 0 \\
    \left[\pi_n, \pi_m\right]
    = \left[\int_{-\pi}^{+\pi} \frac{dk}{2\pi}e^{ink}\pi(k), \int_{-\pi}^{+\pi} \frac{dp}{2\pi}e^{imp}\pi(p)\right]
    = \int_{-\pi}^{+\pi} \frac{dk}{2\pi} \int_{-\pi}^{+\pi}  \frac{dp}{2\pi} e^{i(nk + mp)} \left[\pi(k), \pi(p)\right] = 0 \\
    \left[\phi_n, \pi_m\right]
    = \left[\int_{-\pi}^{+\pi} \frac{dk}{2\pi}e^{ink}\phi(k), \int_{-\pi}^{+\pi} \frac{dp}{2\pi}e^{imp}\pi(p)\right]
    = \int_{-\pi}^{+\pi} \frac{dk}{2\pi} \int_{-\pi}^{+\pi}  \frac{dp}{2\pi} e^{i(nk + mp)} \left[\phi(k), \pi(p)\right] = \\
    = \int_{-\pi}^{+\pi} \frac{dk}{2\pi} \int_{-\pi}^{+\pi}  \frac{dp}{2\pi} e^{i(nk + mp)} i\sum_{l=-\infty}^{\infty}e^{-il(k + p)} = \\
    = i\sum_{l=-\infty}^{\infty} \int_{-\pi}^{+\pi} \frac{dk}{2\pi} \int_{-\pi}^{+\pi}  \frac{dp}{2\pi} e^{i(nk + mp) - il(k + p)} = \\
    = i\sum_{l=-\infty}^{\infty} \int_{-\pi}^{+\pi} \frac{dk}{2\pi} \int_{-\pi}^{+\pi}  \frac{dp}{2\pi} e^{i(n - l)k}e^{i(m - l)p} = \\
    = i\sum_{l=-\infty}^{\infty} \delta_{nl} \delta_{ml} = i\delta_{nm}  
\end{gathered}
\end{equation}
\subsection{The Hamiltonian based on the Fourier modes}
\begin{equation}
\begin{gathered}
    H = \frac{1}{2} \sum_{n=-\infty}^{\infty}\left[\pi_n^2 + \left(\phi_{n+1} - \phi_n\right)^2 + m^2\phi_n^2\right] = \\
    = \frac{1}{2}
    \sum_{n=-\infty}^{\infty}\left[
        \left(\int_{-\pi}^{+\pi} \frac{dk}{2\pi}e^{ink}\pi(k)\right)^2 + \left(
            \int_{-\pi}^{+\pi} \frac{dk}{2\pi}e^{i(n + 1)k}\phi(k) -
            \int_{-\pi}^{+\pi} \frac{dk}{2\pi}e^{ink}\phi(k)
        \right)^2 + m^2\left(\int_{-\pi}^{+\pi} \frac{dk}{2\pi}e^{ink}\phi(k)\right)^2
    \right]
\end{gathered}
\end{equation}
The momentum dependent term of the Hamiltonian can be simplified using the substitution $p \rightarrow -p$ in the second integral, inverting the integration order and using the relations $\pi(-p) = \pi^\dagger(p)$ and $\sum_n e^{i(k-p)} = 2\pi\delta(k-p)$:
\begin{equation}
\begin{gathered}
    H_p = \frac{1}{2}
    \sum_{n=-\infty}^{\infty}
        \left(\int_{-\pi}^{+\pi} \frac{dk}{2\pi}e^{ink}\pi(k)\right)^2
    = \frac{1}{2}
    \sum_{n=-\infty}^{\infty} 
        \left(\int_{-\pi}^{+\pi} \frac{dp}{2\pi}e^{ink}\pi(k)\right)
        \left(\int_{-\pi}^{+\pi} \frac{dk}{2\pi}e^{inp}\pi(p)\right) = \\
    = \frac{1}{2}
    \sum_{n=-\infty}^{\infty}
        \left(\int_{+\pi}^{-\pi} \frac{-dk}{2\pi}e^{-ink}\pi(-k)\right) 
        \left(\int_{-\pi}^{+\pi} \frac{dp}{2\pi}e^{inp}\pi(p)\right) = \\
    = \frac{1}{2}
    \sum_{n=-\infty}^{\infty}
        \left(\int_{-\pi}^{+\pi} \frac{dk}{2\pi}e^{-ink}\pi^\dagger(k)\right) 
        \left(\int_{-\pi}^{+\pi} \frac{dp}{2\pi}e^{inp}\pi(p)\right) = \\
    = \frac{1}{2}
    \sum_{n=-\infty}^{\infty}
    \int_{-\pi}^{+\pi} \frac{dk}{2\pi}
    \int_{-\pi}^{+\pi} \frac{dp}{2\pi}
    e^{-in(k - p)}\pi^\dagger(k)\pi(p) = \\
    = \frac{1}{2}
    \int_{-\pi}^{+\pi} \frac{dk}{2\pi}
    \int_{-\pi}^{+\pi} \frac{dp}{2\pi}
    \pi^\dagger(k)\pi(p)
    \sum_{n=-\infty}^{\infty}e^{-in(k - p)} = \\
    = \frac{1}{2}
    \int_{-\pi}^{+\pi} \frac{dk}{2\pi}
    \int_{-\pi}^{+\pi} \frac{dp}{2\pi}
    \pi^\dagger(k)\pi(p) 2\pi \delta(k-p)
    = \frac{1}{2}
    \int_{-\pi}^{+\pi} \frac{dk}{2\pi}
    \pi^\dagger(k)\pi(k)
\end{gathered}
\end{equation}
The term of the restoring force follows the exact same logic:
\begin{equation}
\begin{gathered}
    \frac{1}{2}
    \sum_{n=-\infty}^{\infty}
    m^2\left(\int_{-\pi}^{+\pi} \frac{dk}{2\pi}e^{ink}\phi(k)\right)^2
    = \frac{1}{2}
    \int_{-\pi}^{+\pi} \frac{dk}{2\pi}
    m^2\phi^\dagger(k)\phi(k)
\end{gathered}
\end{equation}
And the interaction term is
\begin{equation}
\begin{gathered}
    \frac{1}{2}
    \sum_{n=-\infty}^{\infty}
    \left(
        \int_{-\pi}^{+\pi} \frac{dk}{2\pi}e^{i(n + 1)k}\phi(k) -
        \int_{-\pi}^{+\pi} \frac{dk}{2\pi}e^{ink}\phi(k)
    \right)^2 = \\
    = \frac{1}{2}
    \sum_{n=-\infty}^{\infty}
    \left(
        \int_{-\pi}^{+\pi} \frac{dk}{2\pi}e^{ink}\phi(k) (e^{ik} - 1)
    \right)^2 = \\
    = \frac{1}{2}
    \sum_{n=-\infty}^{\infty}
    \left(
        \int_{-\pi}^{+\pi} \frac{dk}{2\pi}e^{ink}\phi(k) (e^{ik} - 1)
    \right)
    \left(
        \int_{-\pi}^{+\pi} \frac{dp}{2\pi}e^{inp}\phi(p) (e^{ip} - 1)
    \right) = \\
    = \frac{1}{2}
    \sum_{n=-\infty}^{\infty}
        \int_{-\pi}^{+\pi} \frac{dk}{2\pi}\int_{-\pi}^{+\pi} \frac{dp}{2\pi}
        e^{-in(k - p)}\phi^\dagger(k)\phi(p)
        (e^{-ik} - 1)(e^{ip} - 1)
    = \\
    = \frac{1}{2}
        \int_{-\pi}^{+\pi} \frac{dk}{2\pi}
        \phi^\dagger(k)\phi(k)
        (e^{-ik} - 1)(e^{ik} - 1)
    = \\
    = \frac{1}{2}
        \int_{-\pi}^{+\pi} \frac{dk}{2\pi}
        \phi^\dagger(k)\phi(k)
        (2 - 2\cos(k))
    = \\
\end{gathered}
\end{equation}
Putting it all together we get
\begin{equation}
\begin{gathered}
    H = \frac{1}{2}
    \int_{-\pi}^{+\pi} \frac{dk}{2\pi}
    \left[
        \pi^\dagger(k)\pi(k) + 
        \left(
            m^2 + 2(1-\cos(k))
        \right)
        \phi^\dagger(k)\phi(k)
    \right] = \\
    = \frac{1}{2}
    \int_{-\pi}^{+\pi} \frac{dk}{2\pi}
    \left[
        \pi^\dagger(k)\pi(k) + 
        \omega_k^2
        \phi^\dagger(k)\phi(k)
    \right]
\end{gathered}
\end{equation}
With the dispersion relation
\begin{equation}
    \omega_k = \sqrt{m^2 + 2(1-\cos(k))}
\end{equation}
With the lattice spacing $\Delta$, it becomes
\begin{equation}
    \omega_k = \sqrt{m^2 + \frac{2}{\Delta^2}(1-\cos(k\Delta))}
\end{equation}
And at the limit $k\Delta << 1$ we get
\begin{equation}
    \frac{2}{\Delta^2}(1-\cos(k\Delta)) =
    \frac{2}{\Delta^2}\left[1 - (1 - \frac{(k\Delta)^2}{2})\right] =
    \frac{(k\Delta)^2}{\Delta^2} = k^2
\end{equation}
Which is the Klein-Gordon dispersion relation from Special Relatvity:
\begin{equation}
    \omega_k = \sqrt{m^2 + k^2}
\end{equation}
\subsection{Ladder Operators}
\begin{equation}
    a_k = \frac{1}{\sqrt{4\pi \omega_k}}\left(\omega_k\phi(k) + i\pi(k)\right), \qquad
    a_k^\dagger = \frac{1}{\sqrt{4\pi \omega_k}}\left(\omega_k\phi^\dagger(k) - i\pi^\dagger(k)\right)
\end{equation}
The commutation relations are
\begin{equation}
\begin{gathered}
    \left[a_k, a_p\right] = \frac{1}{4\pi \sqrt{\omega_k\omega_p}}
    \left[
        \omega_k\phi(k) + i\pi(k),
        \omega_p\phi(p) + i\pi(p)
    \right] = \\
    = \frac{1}{4\pi \sqrt{\omega_k\omega_p}}
    \left[
    \omega_k\omega_p [\phi(k), \phi(p)]
    + i \left(
        \omega_k[\phi(k),\pi(p)] + \omega_p[\pi(k), \phi(p)]
    \right)
    - [\pi(k), \pi(p)]
    \right] = \\
    = - \frac{1}{4\pi \sqrt{\omega_k\omega_p}}
    \left(
        \omega_k \sum_{n=-\infty}^\infty e^{-in(k+p)} -
        \omega_p \sum_{n=-\infty}^\infty e^{-in(k+p)}
    \right) = \\
    = - \frac{\omega_k - \omega_p}{4\pi \sqrt{\omega_k\omega_p}} 2\pi \delta(k+p) = 0
\end{gathered}
\end{equation}
Where in the last part the delta function forced $k = -p$, in which case $\omega_k = \omega_p$ since $\omega$ is invaraint to the sign of $k$. From this commutation relations the commutation relations $\left[a^\dagger_k, a^\dagger_p\right]$ are obvious to develop using the relations $\phi^\dagger(k) = \phi(-k)$ and $\pi^\dagger(k) = \pi(-k)$, and a variable substitution back $k \rightarrow -k'$ and $p \rightarrow -p'$:
\begin{equation}
\begin{gathered}
    \left[a_k^\dagger, a_p^\dagger\right] = \frac{1}{4\pi \sqrt{\omega_k\omega_p}}
    \left[
        \omega_k\phi^\dagger(k) - i\pi^\dagger(k),
        \omega_p\phi^\dagger(p) - i\pi^\dagger(p)
    \right] = \\
     = \frac{1}{4\pi \sqrt{\omega_k\omega_p}}
    \left[
        \omega_k\phi(-k) - i\pi(-k),
        \omega_p\phi(-p) - i\pi(-p)
    \right] = \\
     = \frac{1}{4\pi \sqrt{\omega_{-k'}\omega_{-p'}}}
    \left[
        \omega_{-k'}\phi(k') - i\pi(k'),
        \omega_{-p'}\phi(p') - i\pi(p')
    \right] = \\
    = \frac{\omega_{-k'} - \omega_{-p'}}{4\pi \sqrt{\omega_{-k'}\omega_{-p'}}} 2\pi \delta(k'+p') = 0
\end{gathered}
\end{equation}
Where we again used the fact that $\omega_k = \omega_{-k}$, this time to show that if $p' = k'$ then $\omega_{-p'} = \omega_{p'} = \omega_{k}$. The last commutation relation we want to check is $\left[a_k, a^\dagger_p\right]$:
\begin{equation}
\begin{gathered}
    \left[a_k, a_p^\dagger\right] = \frac{1}{4\pi \sqrt{\omega_k\omega_p}}
    \left[
        \omega_k\phi(k) + i\pi(k),
        \omega_p\phi^\dagger(p) - i\pi^\dagger(p)
    \right] = \\
     = \frac{1}{4\pi \sqrt{\omega_k\omega_p}}
    \left[
        \omega_k\phi(k) + i\pi(k),
        \omega_p\phi(-p) - i\pi(-p)
    \right] = \\
     = \frac{1}{4\pi \sqrt{\omega_k\omega_{-p'}}}
    \left[
        \omega_k\phi(k) + i\pi(k),
        \omega_{-p'}\phi(p') - i\pi(p')
    \right] = \\
     = \frac{1}{4\pi \sqrt{\omega_k\omega_{-p'}}}
    \left(
        - i\omega_k\left[\phi(k),\pi(p')\right]
        + i\omega_{-p'}\left[\pi(k),\phi(p')\right]
    \right) = \\
     = \frac{\omega_k + \omega_{-p'}}{4\pi \sqrt{\omega_k\omega_{-p'}}}
    2\pi \delta(k+p') = \delta(k + p') = \delta(k - p)
\end{gathered}
\end{equation}
Where in the last line, the only case when the delta function is not zero is the case where the fraction evaluates to one.
\newline
Now we will develop the Hamiltonian in terms of $a(k)$ and $a^\dagger(k)$.
To do that we will add and subtract from the integrand in equation 24 the term $i\omega_k\pi(k)\phi^\dagger(k)$. Then, we will take the subtracted term and factorize only it by changing the integration variable $k \rightarrow -k$. This only moves the dagger to the $\pi(k)$ field, since $\pi^\dagger(k) = \pi(-k)$ and $\phi^\dagger(k) = \phi(-k)$, and the change of sign in the differential cancels out with the change in order of the integration limits:
\begin{equation}
\begin{gathered}
    H = \frac{1}{2}
    \int_{-\pi}^{+\pi} \frac{dk}{2\pi}
    \left[
        \pi^\dagger(k)\pi(k) + 
        \omega_k^2
        \phi^\dagger(k)\phi(k)
    \right] = \\
    = \int_{-\pi}^{+\pi} \frac{dk}{4\pi}
    \left[
        \omega_k^2
        \phi^\dagger(k)\phi(k) +
        \pi^\dagger(k)\pi(k) + 
        i\omega_k\left(\pi(k)\phi^\dagger(k) - \pi(k)\phi^\dagger(k)\right)
    \right] = \\
    = \int_{-\pi}^{+\pi} \frac{dk}{4\pi}
    \left[
        \omega_k^2
        \phi^\dagger(k)\phi(k) +
        \pi^\dagger(k)\pi(k) + 
        i\omega_k\left(\pi(k)\phi^\dagger(k) - \pi^\dagger(k)\phi(k)\right)
    \right] = \\
    = \int_{-\pi}^{+\pi} \frac{dk}{4\pi}
    \left[
        \omega_k^2
        \phi^\dagger(k)\phi(k) +
        \pi^\dagger(k)\pi(k) + 
        i\omega_k\left(\phi^\dagger(k)\pi(k) - \pi^\dagger(k)\phi(k) - \left[\phi^\dagger(k), \pi(k)\right]\right)
    \right] = \\
    = \int_{-\pi}^{+\pi} \frac{dk}{4\pi}
    \left[
        \omega_k^2
        \phi^\dagger(k)\phi(k) +
        \pi^\dagger(k)\pi(k) + 
        i\omega_k\left(\phi^\dagger(k)\pi(k) - \pi^\dagger(k)\phi(k) - \left[\phi(-k), \pi(k)\right]\right)
    \right] = \\
    = \int_{-\pi}^{+\pi} \frac{dk}{4\pi}
    \left[
        \omega_k^2
        \phi^\dagger(k)\phi(k) +
        \pi^\dagger(k)\pi(k) + 
        i\omega_k\left(\phi^\dagger(k)\pi(k) - \pi^\dagger(k)\phi(k) - i2\pi\delta(-k+k)\right)
    \right] = \\
    = \int_{-\pi}^{+\pi} \frac{dk}{4\pi}
    \left[
        \omega_k^2
        \phi^\dagger(k)\phi(k) +
        i\omega_k\phi^\dagger(k)\pi(k) - i\omega_k\pi^\dagger(k)\phi(k) + 
        \pi^\dagger(k)\pi(k) + 2\pi\omega_k\delta(0)
    \right]
\end{gathered}
\end{equation}
The first four terms can be factorized:
\begin{equation}
\begin{gathered}
    \int_{-\pi}^{+\pi} \frac{dk}{4\pi}
    \left[
        \left(\omega_k\phi^\dagger(k) - i\pi^\dagger(k)\right)
        \left(\omega_k\phi(k) + i\pi(k)\right)
        +2\pi\omega_k\delta(0)   
    \right] = \\
    = \int_{-\pi}^{+\pi}  dk \omega_k
    \left[
        \frac{1}{4\pi\omega_k}\left(\omega_k\phi^\dagger(k) - i\pi^\dagger(k)\right)
        \left(\omega_k\phi(k) + i\pi(k)\right)
        + \frac{1}{2}\delta(0)   
    \right] = \\
    = \int_{-\pi}^{+\pi}  dk \omega_k
    \left[a^\dagger(k)a(k)+ \frac{1}{2}\delta(0)   
    \right]
\end{gathered}
\end{equation}
\subsection{The Vaccum Energy relation to N}
In this situation we switch the formula for the commutator $\left[\phi(k),\pi(p)\right]$:
\begin{equation}
    \left[\phi(k),\pi(p)\right] = \sum_{-N}^{N} e^{-in(k+p)}
\end{equation}
And for the situation of the vaccum energy, we got
\begin{equation}
    \left[\phi(-k),\pi(k)\right] = \sum_{-N}^{N} e^{-in(-k+k)} = \sum_{-N}^{N} 1 = 2N+1
\end{equation}
Meaning that the energy is linear to the number of points.
\end{document}